\documentclass[12pt,a4paper]{article}
\usepackage[utf8]{inputenc}
\usepackage[T1]{fontenc}
\usepackage{mathptmx}
\usepackage[margin=2cm]{geometry}
\usepackage{hyperref}
\hypersetup{colorlinks=false, pdfborder={0 0 0}}
\usepackage{titling}
\setlength{\droptitle}{-2.5cm}
\pretitle{\begin{center}\Large}
\posttitle{\par\end{center}\vspace{-0.5em}}
\postauthor{\end{tabular}\par\end{center}\vspace{-0.5em}}
\postdate{\par\end{center}\vspace{-1.5em}}

\title{\textbf{Threshold-Aware Similarity Caching for\\
Named Data Networks:\\
Bounding Retry Delay under Congestion using ns-3}}
\author{Somyak Priyadarshi Mohanta (241CS257)\\
Pranav Shaji (241CS243)\\
Chris Tony (241CS222)\\
\textbf{Department of Computer Science and Engineering}\\
\textbf{National Institute of Technology Karnataka, Surathkal}}
\date{\today}

\begin{document}
\maketitle

\section*{Abstract}
Named Data Networking (NDN) lets every router cache the Data packets it forwards, but, like any
cache, it answers a request only with the exact object that was named \cite{ndn}. Similarity
caching relaxes this rule: on a miss, a cache may return the most similar object it holds, which
suits image retrieval, recommendation and semantic caching of large language model queries
\cite{neglia}, \cite{vcache}. The user accepts the answer only if it is similar enough, and
otherwise requests again. Making this work in a multi-hop network is difficult for four reasons: a
cache does not know the similarity the user will accept, so it may send an answer certain to be
rejected; NDN delivers Data only under the name of the pending Interest, so a substitute object
cannot be returned as is; cache contents do not change within a retry interval, so a rejected user
keeps getting the same answer; and the value of a nearby answer depends on upstream queueing and
loss. Nakamura and Kamiyama \cite{nakamura} model the similar-content delivery delay on general
cache networks, but assume a fixed answer probability $\alpha$, lossless links of infinite
bandwidth and retries that redraw cache contents independently, with numerical results only.
Keeping contents fixed across retries on their topology, we find their model underestimates the
mean delay up to 2.1 times at $\alpha = 0.5$ and up to 1700 times at $\alpha = 0.9$; other studies
treat a single cache \cite{neglia} or offline placement \cite{garetto}, not packet-level retries.
We therefore
propose threshold-aware similarity caching for NDN, in which Interests carry the user's minimum
acceptable similarity and caches answer or forward based on the measured upstream delay. The
objectives are: (i) to implement similarity caching with consumer retries in ns-3 using ndnSIM
\cite{ndnsim}, reproduce the base paper on its tandem topology and measure the delay when cache
contents persist across retries; (ii) to design threshold-carrying Interests with bounded retries
and derive their delay in closed form; and (iii) to design a congestion- and loss-aware
answer-or-forward rule and compare it with fixed-$\alpha$ caching under bottleneck links and burst
loss on mean and 95th-percentile delay, delivered similarity and upstream traffic.

\vspace{-0.8em}
{\small
\begin{thebibliography}{9}
\setlength{\itemsep}{1pt}

\bibitem{ndn}
L. Zhang \textit{et al.}, ``Named data networking,'' \textit{ACM SIGCOMM Comput. Commun. Rev.},
vol.~44, no.~3, pp.~66--73, 2014.

\bibitem{neglia}
G. Neglia, M. Garetto, and E. Leonardi, ``Similarity caching: Theory and algorithms,''
\textit{IEEE/ACM Trans. Netw.}, vol.~30, no.~2, pp.~475--486, 2022.

\bibitem{vcache}
L. G. Schroeder \textit{et al.}, ``vCache: Verified semantic prompt caching,'' arXiv:2502.03771,
2025.

\bibitem{nakamura}
R. Nakamura and N. Kamiyama, ``Analysis of similarity caching on general cache networks,''
\textit{IEEE Access}, vol.~12, pp.~163338--163348, 2024.

\bibitem{garetto}
M. Garetto, E. Leonardi, and G. Neglia, ``Content placement in networks of similarity caches,''
\textit{Comput. Netw.}, vol.~201, Art. no.~108570, 2021.

\bibitem{ndnsim}
S. Mastorakis, A. Afanasyev, and L. Zhang, ``On the evolution of ndnSIM: An open-source simulator
for NDN experimentation,'' \textit{ACM SIGCOMM Comput. Commun. Rev.}, vol.~47, no.~3,
pp.~19--33, 2017.

\end{thebibliography}
}

\end{document}
